\documentclass[12pt]{article}
\usepackage{amsmath}
\usepackage{amsfonts}
\usepackage{float}
\usepackage{fancyhdr}
\usepackage{graphicx}
\usepackage{setspace}
\usepackage[colorlinks=true,linkcolor=blue, citecolor=red]{hyperref}
\usepackage{url}
\usepackage[top=.75in, left=.75in, right=.75in, bottom=1in]{geometry}

% For algorithm
\usepackage{algorithm}
\usepackage{algpseudocode}

% ============ CODE ============
\usepackage{listings}
\usepackage{xcolor}
\definecolor{codegreen}{rgb}{0,0.6,0}
\definecolor{codegray}{rgb}{0.5,0.5,0.5}
\definecolor{codepurple}{rgb}{0.58,0,0.82}
\definecolor{backcolour}{rgb}{0.95,0.95,0.92}

% Styling for the code.
\lstdefinestyle{mystyle}{
    backgroundcolor=\color{backcolour},   
    commentstyle=\color{codegreen},
    keywordstyle=\color{magenta},
    numberstyle=\tiny\color{codegray},
    stringstyle=\color{codepurple},
    basicstyle=\ttfamily\footnotesize,
    breakatwhitespace=false,         
    breaklines=true,                 
    captionpos=b,                    
    keepspaces=true,                 
    numbers=left,                    
    numbersep=5pt,                  
    showspaces=false,                
    showstringspaces=false,
    showtabs=false,                  
    tabsize=2
}
\lstset{style=mystyle}

% Table packages and Unicode / Font support
\usepackage{array}
\usepackage{longtable}
\usepackage{enumitem}
\usepackage{booktabs}
\usepackage[utf8]{inputenc}
\DeclareUnicodeCharacter{2264}{\ensuremath{\le}}
\DeclareUnicodeCharacter{2194}{\ensuremath{\leftrightarrow}}
\usepackage[T1]{fontenc}
\usepackage{lmodern}
\graphicspath{{../../../../templates/}{../../prototype/}{figures/}{./}}
\setlength{\emergencystretch}{3em}
\setlist[enumerate]{label=\arabic*.,leftmargin=*}
\setlist[itemize]{leftmargin=*}

% ============ HEADER AND FOOTER ============
\setlength{\headheight}{29.43912pt}

\pagestyle{fancy}
\fancyfoot{}
\fancyfoot[R]{Page \thepage}

\setlength{\parindent}{0pt}
\renewcommand{\baselinestretch}{1.5}

\newcommand{\coursename}{Introduction of Software Engineering}
\newcommand{\reportname}{Studitopia}
\newcommand{\reporttitle}{\href{https://github.com/RoundFloor/Studitopia.git}{Studitopia}}

\newcommand{\studentname}{%
Nguyen Thanh An (24125024)\\
Nguyen Chi Bao (24125025)\\
Tran Nguyen Khanh Binh (24125026)\\
Tran Canh Anh Tuan (24125048)\\
Nguyen Le Quoc Huy (24125100)
}
\newcommand{\teachername}{Ho Tuan Thanh, Nguyen Thi Minh Tuyen}

% Header
\lhead{\reporttitle}
\rhead{
University of Science - VNUHCM\\
\coursename
}

% Footer
\newcommand{\leftfooter}{}
\lfoot{\leftfooter}

% ============ DOCUMENT ============
\begin{document}

\pagenumbering{roman}
\begin{titlepage}
\newcommand{\HRule}{\rule{\linewidth}{0.5mm}}
\centering

\textsc{\LARGE VNU-HCM}\\[0.5cm]
\textsc{\Large University of Science}\\[0.5cm]
\textsc{\large Faculty of Information Technology}\\[0.5cm]
\textsc{Computer Science}\\[0.5cm]

\HRule \\[0.4cm]
{ 
\huge{\bfseries{\reporttitle}}\\[0.5cm]
\large{\bfseries{UC-01 --- Register Account and Verify University Email}}
}\\[0.4cm]
\HRule \\[0.5cm]

\textbf{\large Subject: \coursename}\\[0.5cm]

\begin{minipage}[t]{0.55\textwidth}
\begin{flushleft} \large
\emph{Students:}\\
\studentname
\end{flushleft}
\end{minipage}
~
\begin{minipage}[t]{0.4\textwidth}
\begin{flushright} \large
\emph{Supervising Teachers:} \\
\teachername
\end{flushright}
\end{minipage}\\[1cm]

{\large \today}\\[1cm]

\includegraphics[scale=.20]{img/hcmus-logo.png}\\[1cm] 

\vfill
\end{titlepage}	

\tableofcontents
\pagebreak

\pagenumbering{arabic}
\setcounter{page}{1}

\hypertarget{sec:uc-01-register-account-and-verify-university-email}{}
\addcontentsline{toc}{section}{UC-01 --- Register Account and Verify University Email}
\section*{UC-01 --- Register Account and Verify University Email}
\par\medskip

\textbf{Version:} \texttt{2.0}
\par\medskip

\textbf{Performed by:} [Nguyen Thanh An] | \textbf{Edited by:} [Nguyen Thanh An]
\par\medskip

\begingroup
\small
\renewcommand{\baselinestretch}{1.0}\selectfont
\renewcommand{\arraystretch}{1.15}
\setlength{\tabcolsep}{6pt}
\begin{longtable}{|>{\raggedright\arraybackslash}p{0.09\textwidth}|>{\raggedright\arraybackslash}p{0.84\textwidth}|}
\hline
\textbf{Field} & \textbf{Value} \\ \hline
\endhead
Use Case ID & UC-01 \\ \hline
Name & Register Account and Verify University Email \\ \hline
Primary Actor & New University Student \\ \hline
Supporting Actors & Email Delivery Service; Google OAuth 2.0 \\ \hline
Feature Group & Authentication \\ \hline
Priority & High \\ \hline
Related User Stories & US-01, US-07 \\ \hline
Related Requirements & FR-01.1 to FR-01.45; FR-07.1 to FR-07.6; NFR-SEC-2, NFR-SEC-3, NFR-SEC-8, NFR-USE-1 \\ \hline
Source Documents & \texttt{proposals/use\_cases\_model/UC-01-register-account.md}; \texttt{proposals/use\_cases\_description.md} (UC-01); \texttt{proposals/use\_cases\_model.md}; \texttt{proposals/use\_case\_diagrams.md}; \texttt{PA3-Group05/use\_case.md} (section 1); \texttt{proposals/overview.md}; \texttt{proposals/functional\_requirement.md}; \texttt{proposals/non\_functional\_requirement.md}; \texttt{proposals/userstory.md}; \texttt{proposals/feature\_flow.md}; \texttt{proposals/screen\_flow.md}; \texttt{proposals/API.md}; \texttt{proposals/response\_code.md}; \texttt{proposals/settings.md} \\ \hline
\end{longtable}
\endgroup
\par\medskip

\hypertarget{sec:part-a-use-case-model}{}
\addcontentsline{toc}{subsection}{Part A --- Use Case Model}
\subsection*{Part A --- Use Case Model}
\par\medskip

\hypertarget{sec:a1-focused-use-case-diagram}{}
\addcontentsline{toc}{subsubsection}{A1. Focused Use Case Diagram}
\subsubsection*{A1. Focused Use Case Diagram}
\par\medskip

\begin{center}
\includegraphics[width=0.85\textwidth,height=0.75\textheight,keepaspectratio]{"figures/UC-01-register-account_1.png"}
\end{center}
\par\medskip

\textbf{Diagram notation.} Person shape = human actor; cloud = supporting system; blue = selected base use case; purple = included; green = extending; orange = specialized; solid line = association; dashed arrow = include or extend; labelled solid arrow = generalization.
\par\medskip

\hypertarget{sec:a2-actor-register}{}
\addcontentsline{toc}{subsubsection}{A2. Actor Register}
\subsubsection*{A2. Actor Register}
\par\medskip

\begingroup
\small
\renewcommand{\baselinestretch}{1.0}\selectfont
\renewcommand{\arraystretch}{1.15}
\setlength{\tabcolsep}{6pt}
\begin{longtable}{|>{\raggedright\arraybackslash}p{0.14\textwidth}|>{\raggedright\arraybackslash}p{0.12\textwidth}|>{\raggedright\arraybackslash}p{0.66\textwidth}|}
\hline
\textbf{Actor} & \textbf{Type} & \textbf{Description} \\ \hline
\endhead
New University Student & Primary human actor & A prospective user who creates an account and verifies a university email address. \\ \hline
University Student & General human actor & Parent actor of New University Student; specialized actors inherit the general student associations. \\ \hline
Email Delivery Service & Supporting system & Delivers time-limited OTP codes for account verification. \\ \hline
Google OAuth 2.0 & Supporting system & Provides the optional external authentication and registration path. \\ \hline
\end{longtable}
\endgroup
\par\medskip

\hypertarget{sec:a3-relationships}{}
\addcontentsline{toc}{subsubsection}{A3. Relationships}
\subsubsection*{A3. Relationships}
\par\medskip

\hypertarget{sec:a3-1-associations}{}
\paragraph*{A3.1 Associations}
\par\medskip

\begingroup
\small
\renewcommand{\baselinestretch}{1.0}\selectfont
\renewcommand{\arraystretch}{1.15}
\setlength{\tabcolsep}{6pt}
\begin{longtable}{|>{\raggedright\arraybackslash}p{0.15\textwidth}|>{\raggedright\arraybackslash}p{0.34\textwidth}|>{\raggedright\arraybackslash}p{0.10\textwidth}|>{\raggedright\arraybackslash}p{0.33\textwidth}|}
\hline
\textbf{Actor} & \textbf{Use Case} & \textbf{Type} & \textbf{Description} \\ \hline
\endhead
New University Student & UC-01 Register Account and Verify University Email & Primary & Starts registration and email verification. \\ \hline
Google OAuth 2.0 & G-02 Register with Google & Supporting & Authenticates the Google registration path. \\ \hline
Email Delivery Service & I-02 Verify Email OTP & Supporting & Delivers the OTP to the university email address. \\ \hline
\end{longtable}
\endgroup
\par\medskip

\hypertarget{sec:a3-2-generalization-and-specialization}{}
\paragraph*{A3.2 Generalization and Specialization}
\par\medskip

\begingroup
\small
\renewcommand{\baselinestretch}{1.0}\selectfont
\renewcommand{\arraystretch}{1.15}
\setlength{\tabcolsep}{6pt}
\begin{longtable}{|>{\raggedright\arraybackslash}p{0.16\textwidth}|>{\raggedright\arraybackslash}p{0.23\textwidth}|>{\raggedright\arraybackslash}p{0.11\textwidth}|>{\raggedright\arraybackslash}p{0.41\textwidth}|}
\hline
\textbf{Child} & \textbf{Parent} & \textbf{Type} & \textbf{Inherited Behavior or Association} \\ \hline
\endhead
New University Student & University Student & Actor generalization & Inherits the general student associations and adds UC-01. \\ \hline
G-01 Register with University Email & UC-01 Register Account and Verify University Email & Use case generalization & Inherits the UC-01 actor association and I-01 include; adds the mandatory OTP include. \\ \hline
G-02 Register with Google & UC-01 Register Account and Verify University Email & Use case generalization & Inherits the UC-01 actor association and I-01 include; uses Google OAuth instead of OTP. \\ \hline
\end{longtable}
\endgroup
\par\medskip

\hypertarget{sec:a3-3-include}{}
\paragraph*{A3.3 Include}
\par\medskip

\begingroup
\small
\renewcommand{\baselinestretch}{1.0}\selectfont
\renewcommand{\arraystretch}{1.15}
\setlength{\tabcolsep}{6pt}
\begin{longtable}{|>{\raggedright\arraybackslash}p{0.26\textwidth}|>{\raggedright\arraybackslash}p{0.16\textwidth}|>{\raggedright\arraybackslash}p{0.50\textwidth}|}
\hline
\textbf{Base Use Case} & \textbf{Included Use Case} & \textbf{Meaning} \\ \hline
\endhead
UC-01 Register Account and Verify University Email & I-01 Validate Registration Data & Validate required fields, username, email, password, confirmation, CAPTCHA, and uniqueness rules. \\ \hline
G-01 Register with University Email & I-02 Verify Email OTP & Send and verify a 6-digit OTP with a 5-minute expiry. \\ \hline
\end{longtable}
\endgroup
\par\medskip

\hypertarget{sec:a3-4-extend}{}
\paragraph*{A3.4 Extend}
\par\medskip

\begingroup
\small
\renewcommand{\baselinestretch}{1.0}\selectfont
\renewcommand{\arraystretch}{1.15}
\setlength{\tabcolsep}{6pt}
\begin{longtable}{|>{\raggedright\arraybackslash}p{0.10\textwidth}|>{\raggedright\arraybackslash}p{0.11\textwidth}|>{\raggedright\arraybackslash}p{0.25\textwidth}|>{\raggedright\arraybackslash}p{0.45\textwidth}|}
\hline
\textbf{Extending Use Case} & \textbf{Base Use Case} & \textbf{Condition} & \textbf{Added Behavior} \\ \hline
\endhead
E-01 Resend OTP & I-02 Verify Email OTP & The user requests a new OTP before verification. & Invalidate the previous OTP, generate a new code, and restart the 5-minute expiry period. \\ \hline
\end{longtable}
\endgroup
\par\medskip

\hypertarget{sec:a4-model-notes}{}
\addcontentsline{toc}{subsubsection}{A4. Model Notes}
\subsubsection*{A4. Model Notes}
\par\medskip

\begin{enumerate}
\item UC-01 is the selected general registration use case. G-01 and G-02 model the two registration paths without duplicating the base use case.
\item OTP verification is modeled under G-01 because it is the mandatory verification step for university-email registration. G-02 uses Google OAuth instead.
\item G-01 and G-02 inherit the UC-01 actor association and the I-01 validation behavior through use case generalization.
\end{enumerate}
\par\medskip

\hypertarget{sec:part-b-use-case-specification}{}
\addcontentsline{toc}{subsection}{Part B --- Use Case Specification}
\subsection*{Part B --- Use Case Specification}
\par\medskip

\hypertarget{sec:b1-brief-description}{}
\addcontentsline{toc}{subsubsection}{B1. Brief Description}
\subsubsection*{B1. Brief Description}
\par\medskip

A new university student creates a verified account and gains access to protected Studitopia features. The student can register with a supported university email address and a one-time password, or use Google OAuth. After successful registration and verification, the student can configure learning preferences and access the Home page.
\par\medskip

\hypertarget{sec:b2-preconditions}{}
\addcontentsline{toc}{subsubsection}{B2. Preconditions}
\subsubsection*{B2. Preconditions}
\par\medskip

\begin{enumerate}
\item The user does not have an active authenticated session.
\item The supported university and faculty email-domain catalog is available for university-email registration.
\item The email delivery service is available for the standard registration path.
\item Google OAuth 2.0 is configured for the optional Google registration path.
\end{enumerate}
\par\medskip

\hypertarget{sec:b3-postconditions}{}
\addcontentsline{toc}{subsubsection}{B3. Postconditions}
\subsubsection*{B3. Postconditions}
\par\medskip

\begin{enumerate}
\item \textbf{Success:} A user record exists with a unique username and email address, and the account is activated and email-verified.
\item \textbf{Success:} The user's password is stored only as a secure one-way hash; Google registration associates the authenticated Google identity with the account.
\item \textbf{Success:} The user is authorized to continue to Preference Settings and then the Home page.
\item \textbf{Failure:} No active account is created, and the user receives an error message that allows correction and retry.
\item \textbf{Failure:} An unverified account, if created, cannot access protected application features.
\end{enumerate}
\par\medskip

\hypertarget{sec:b4-main-success-scenario-basic-flow}{}
\addcontentsline{toc}{subsubsection}{B4. Main Success Scenario (Basic Flow)}
\subsubsection*{B4. Main Success Scenario (Basic Flow)}
\par\medskip

\begin{enumerate}
\item The user opens the registration page and selects \textbf{Register}.
\item The system displays Display Name, Username, Email, Password, Confirm Password, CAPTCHA, and a \textbf{Continue with Google} option.
\item The user enters the required details and completes the CAPTCHA.
\item The frontend validates that all required fields are present, the email format is valid, the username contains only alphanumeric characters, and the password satisfies the password policy.
\item The frontend verifies that Confirm Password matches Password.
\item The user submits the registration form.
\item The backend validates the submitted data, checks that the email and username are unique, and verifies that the email domain belongs to a supported university or faculty when university-email registration is used.
\item The backend stores a pending account with a securely hashed password and generates a time-limited OTP.
\item The email delivery service sends the OTP to the registered email address, and the system redirects the user to the OTP verification page.
\item The user enters the 6-digit OTP and submits it.
\item The backend verifies that the OTP is valid, unexpired, unused, and correctly formatted.
\item The system activates the account, updates the email-verification status, and authorizes access to Preference Settings.
\item The user completes or updates learning preferences and continues to the Home page.
\end{enumerate}
\par\medskip

\hypertarget{sec:b5-alternative-flows}{}
\addcontentsline{toc}{subsubsection}{B5. Alternative Flows}
\subsubsection*{B5. Alternative Flows}
\par\medskip

Unless otherwise stated, an alternative flow resumes at the next step of the basic flow.
\par\medskip

\hypertarget{sec:a1-google-registration}{}
\paragraph*{A1. Google registration}
\par\medskip

At basic flow step 2, the user selects \textbf{Continue with Google}. The system initiates Google OAuth, validates the authentication response, and creates or associates an account when registration requirements are satisfied. The user then continues to the appropriate next screen, normally Preference Settings.
\par\medskip

\hypertarget{sec:a2-invalid-client-side-input}{}
\paragraph*{A2. Invalid client-side input}
\par\medskip

At basic flow step 4 or 5, validation fails. The system displays field-level errors, keeps the entered values where appropriate, and does not submit the request until the errors are corrected.
\par\medskip

\hypertarget{sec:a3-duplicate-email-or-username}{}
\paragraph*{A3. Duplicate email or username}
\par\medskip

At basic flow step 7, the backend reports that the email or username is already registered. The system displays a clear error and offers navigation to Login or Password Recovery.
\par\medskip

\hypertarget{sec:a4-captcha-verification-failure}{}
\paragraph*{A4. CAPTCHA verification failure}
\par\medskip

At basic flow step 7, CAPTCHA verification fails. The system rejects the request, shows an error, and allows the user to retry.
\par\medskip

\hypertarget{sec:a5-resend-otp}{}
\paragraph*{A5. Resend OTP}
\par\medskip

At basic flow step 10, the user requests a new OTP. The system generates a new 6-digit OTP, invalidates the previous OTP, sends the new code, and restarts the 5-minute expiry period.
\par\medskip

\hypertarget{sec:a6-user-cancels-registration}{}
\paragraph*{A6. User cancels registration}
\par\medskip

The user leaves the registration flow before activation. No active account is available, and the pending verification data remains subject to the system's expiry and cleanup policy.
\par\medskip

\hypertarget{sec:b6-exception-flows}{}
\addcontentsline{toc}{subsubsection}{B6. Exception Flows}
\subsubsection*{B6. Exception Flows}
\par\medskip

\hypertarget{sec:e1-otp-invalid-expired-or-already-used}{}
\paragraph*{E1. OTP invalid, expired, or already used}
\par\medskip

The system rejects the OTP, displays an appropriate message, and keeps the account unverified. The user may request a new OTP.
\par\medskip

\hypertarget{sec:e2-otp-email-delivery-failure}{}
\paragraph*{E2. OTP email delivery failure}
\par\medskip

The account remains unverified and protected features remain unavailable. The system informs the user and allows another OTP request when permitted.
\par\medskip

\hypertarget{sec:e3-google-authentication-fails-or-is-cancelled}{}
\paragraph*{E3. Google authentication fails or is cancelled}
\par\medskip

No account is created or changed by the failed attempt. The system displays an error and allows the user to retry or use manual registration.
\par\medskip

\hypertarget{sec:e4-database-failure-during-account-creation}{}
\paragraph*{E4. Database failure during account creation}
\par\medskip

The registration transaction is rolled back, a generic error is shown, and the user can retry without creating a partial or duplicate account.
\par\medskip

\hypertarget{sec:e5-unverified-account-attempts-protected-access}{}
\paragraph*{E5. Unverified account attempts protected access}
\par\medskip

The backend denies access to protected features until email verification succeeds.
\par\medskip

\hypertarget{sec:b7-business-rules}{}
\addcontentsline{toc}{subsubsection}{B7. Business Rules}
\subsubsection*{B7. Business Rules}
\par\medskip

\begin{enumerate}
\item Display Name, Username, Email, Password, and Confirm Password are required.
\item A username must be unique and contain only alphanumeric characters.
\item An email address must be unique and valid; university-email registration requires a supported university or faculty domain.
\item A password must contain at least 8 characters, including at least one lowercase letter, one uppercase letter, one number, and one special character.
\item Passwords and Confirm Password fields are hidden by default and support a visibility toggle.
\item The OTP is a 6-digit numeric code, expires after 5 minutes, is single-use, and is invalidated when a new OTP is generated.
\item An account cannot access protected application features before successful email verification.
\item The password is persisted only as a secure one-way hash.
\end{enumerate}
\par\medskip

\hypertarget{sec:b8-special-requirements}{}
\addcontentsline{toc}{subsubsection}{B8. Special Requirements}
\subsubsection*{B8. Special Requirements}
\par\medskip

\begingroup
\small
\renewcommand{\baselinestretch}{1.0}\selectfont
\renewcommand{\arraystretch}{1.15}
\setlength{\tabcolsep}{6pt}
\begin{longtable}{|>{\raggedright\arraybackslash}p{0.10\textwidth}|>{\raggedright\arraybackslash}p{0.83\textwidth}|}
\hline
\textbf{ID} & \textbf{Requirement} \\ \hline
\endhead
NFR-SEC-2 & User credentials must be stored using one-way password hashing with a secure hashing algorithm. \\ \hline
NFR-SEC-3 & Sensitive communication must be encrypted using TLS 1.2 or higher. \\ \hline
NFR-SEC-8 & The system must protect against common web security vulnerabilities, including SQL injection, XSS, and CSRF. \\ \hline
NFR-USE-1 & A new university student must be able to complete registration within 5 minutes without external assistance. \\ \hline
Design constraint & The OTP is 6 digits, single-use, and expires after 5 minutes. The password policy and field visibility rules must be enforced by both frontend and backend validation. \\ \hline
\end{longtable}
\endgroup
\par\medskip

\hypertarget{sec:b9-extension-points}{}
\addcontentsline{toc}{subsubsection}{B9. Extension Points}
\subsubsection*{B9. Extension Points}
\par\medskip

\begingroup
\small
\renewcommand{\baselinestretch}{1.0}\selectfont
\renewcommand{\arraystretch}{1.15}
\setlength{\tabcolsep}{6pt}
\begin{longtable}{|>{\raggedright\arraybackslash}p{0.32\textwidth}|>{\raggedright\arraybackslash}p{0.61\textwidth}|}
\hline
\textbf{Extension Point} & \textbf{Location in Flow of Events} \\ \hline
\endhead
Google registration & Basic flow step 2; alternative flow A1 \\ \hline
Client-side validation recovery & Basic flow steps 4-5; alternative flow A2 \\ \hline
Duplicate-account handling & Basic flow step 7; alternative flow A3 \\ \hline
CAPTCHA retry & Basic flow step 7; alternative flow A4 \\ \hline
OTP resend & Basic flow step 10; alternative flow A5 and exceptions E1-E2 \\ \hline
Registration cancellation & Any step before activation; alternative flow A6 \\ \hline
Unverified-access recovery & Protected-feature access; exception E5 \\ \hline
\end{longtable}
\endgroup
\par\medskip

\hypertarget{sec:b10-interface-and-data-notes}{}
\addcontentsline{toc}{subsubsection}{B10. Interface and Data Notes}
\subsubsection*{B10. Interface and Data Notes}
\par\medskip

\begin{itemize}
\item \textbf{Screen/flow path:} \texttt{Authentication → Register → OTP Verification → Preference Settings → Home} (\texttt{feature\_flow.md}); \texttt{IsLogin = False → Authentication}, with \texttt{Login ↔ Register} and \texttt{Register → OTP → Preference Settings} (\texttt{screen\_flow.md}).
\item \textbf{Registration UI:} Display Name, Username, Email, Password, Confirm Password, CAPTCHA, \textbf{Register} submission, and \textbf{Continue with Google} option; Password and Confirm Password are hidden by default with a visibility toggle.
\item \textbf{OTP UI:} 6-slot numeric OTP input, masked recipient-email notice, \textbf{Resend code} control, expiry hint, and \textbf{Submit Button}; OTP feedback is shown on invalid or expired codes.
\item \textbf{Preference step:} During onboarding the user reaches Preference Settings after OTP verification and continues to Home; \texttt{settings.md} recommends \texttt{/settings/preferences} as the route.
\item \textbf{Documented API:} \texttt{POST /auth/register} with body \texttt{displayname}, \texttt{username}, \texttt{email}, \texttt{password}; \texttt{GET /auth/validate} is marked "Temporarily Disabled" and accepts optional \texttt{username}, \texttt{email}, \texttt{displayname}, \texttt{password} (not recommended) query parameters.
\item \textbf{Register responses:} \texttt{201 SUCCESS} with \texttt{access\_token} (\texttt{token}, \texttt{created}, \texttt{expired}); \texttt{400 BAD\_REQUEST}; \texttt{409 CONFLICT} with \texttt{DUPPLICATED\_USERNAME}, \texttt{DUPLICATED\_EMAIL}, \texttt{UNREGISTED\_EMAIL\_DOMAIN}; \texttt{422 UNPROCESSABLE} with \texttt{INVALID\_DISPLAYNAME}, \texttt{INVALID\_USERNAME}, \texttt{INVALID\_EMAIL}, \texttt{INVALID\_PASSWORD}; \texttt{500 INTERNAL\_SERVER\_ERROR} with \texttt{DATABASE\_ERROR}, \texttt{UNKNOWN\_ERROR}.
\item \textbf{Error codes:} \texttt{0 SUCCESS}, \texttt{2002 UNKNOWN ERROR}, \texttt{45000 UNKNOWN DATABASE ERROR}, \texttt{45001 INVALID\_DISPLAYNAME}, \texttt{45002 INVALID\_EMAIL}, \texttt{45003 INVALID\_USERNAME}, \texttt{45004 INVALID\_PASSWORD}, \texttt{46001 UNREGISTED\_EMAIL\_DOMAIN}, \texttt{47001 DUPLICATED\_EMAIL}, \texttt{47002 DUPLICATED\_USERNAME}.
\item \textbf{Account/status data:} \texttt{studitopia.register(displayname, username, email, password)} returns \texttt{session\_id}, \texttt{user\_id}, \texttt{status}; status values are \texttt{0 (inactive)}, \texttt{1 (preprofile)}, \texttt{2 (active)}, \texttt{3 (suspend)}. Account registered must be inactive; password must be hashed; HMAC-SHA256 is specified for \texttt{refresh\_token} and \texttt{access\_token}.
\item \textbf{Validation constraints:} Required fields; unique alphanumeric username; valid and unique email; supported university or faculty domain; minimum 8-character password with lowercase, uppercase, number, and special character; Confirm Password match; CAPTCHA verification on frontend and backend.
\item \textbf{OTP constraints:} 6-digit numeric code, sent to the registered email, stored as OTP or secure representation with expiration and verification status, expires after 5 minutes, single-use, invalidated on resend.
\item \textbf{Notifications and privacy:} OTP delivery is performed by the Email Delivery Service; sensitive communication uses TLS 1.2 or higher; password hashes are never displayed; unverified accounts are restricted from protected application features.
\end{itemize}
\par\medskip

\hypertarget{sec:b11-source-alignment-notes}{}
\addcontentsline{toc}{subsubsection}{B11. Source Alignment Notes}
\subsubsection*{B11. Source Alignment Notes}
\par\medskip

\begin{itemize}
\item API coverage is partial: no endpoint is documented for OTP submission, OTP resend, Google OAuth registration, or the Preference Settings transition; only \texttt{POST /auth/register} and the temporarily disabled \texttt{GET /auth/validate} are present.
\item \texttt{POST /auth/register} returns an \texttt{access\_token}, while the API also requires the registered account to be inactive and UC-01 postconditions forbid protected access before email verification; the intended token scope is not specified.
\item CAPTCHA verification is required by \texttt{FR-01.17} and business rules, but no dedicated CAPTCHA error code is listed in \texttt{response\_code.md}.
\item Error-name spelling differs across sources: \texttt{API.md} uses \texttt{DUPPLICATED\_EMAIL}, while \texttt{response\_code.md} uses \texttt{DUPLICATED\_EMAIL}; the unusual \texttt{UNREGISTED\_EMAIL\_DOMAIN} spelling is preserved from both.
\item The focused model states G-02 uses Google OAuth instead of OTP, but the sources do not explicitly state whether the Google-authenticated email is treated as verified; A1 only says the user continues, normally to Preference Settings.
\item \texttt{/settings/preferences} is a recommended route in \texttt{settings.md}, not a route defined in the proposal flow documents.
\end{itemize}
\par\medskip

\par\medskip
\hypertarget{sec:ui-prototype-specimens}{}
\addcontentsline{toc}{subsection}{UI Prototype Specimens}
\subsection*{UI Prototype Specimens}
\par\medskip
\begin{figure}[H]
\centering
\includegraphics[width=0.75\textwidth,height=0.42\textheight,keepaspectratio]{"../../prototype/otp_desktop.png"}
\par\vspace{2pt}
{\footnotesize \textbf{UI Prototype:} OTP - Desktop}
\end{figure}
\par\smallskip
\begin{figure}[H]
\centering
\includegraphics[width=0.75\textwidth,height=0.42\textheight,keepaspectratio]{"../../prototype/register_desktop.png"}
\par\vspace{2pt}
{\footnotesize \textbf{UI Prototype:} Register - Desktop}
\end{figure}
\par\smallskip

\end{document}
